%============================= 7. CONCLUSIONS =================================
\isec{\faFlagCheckered}{Conclusions}
\begin{itemize}
  \item \textbf{Mission completed on the real robot.} In the filmed run the EP swept in place with T3 hidden, centered
        it, pursued it along a cubic flat-output plan and stopped 20\,cm before it (tape). The two earlier floor runs
        stopped at 21 and 20\,cm.
  \item \textbf{Flatness gives speed limits by construction.} With $T=\max(4,1.6(d_0-0.20)/v_{max})$ and
        $\dot x(0)=0.06$\,m/s the cubic peaks at 0.1371\,m/s for every $d_0>0.6$\,m; the measured command matched it.
  \item \textbf{Voltage telemetry behaves as the model predicts.} 0.456\,V in a steady search turn, about 1.06\,V at
        cruise, one-tick spikes at command steps, small negative voltages while braking and exactly 0.000\,V after the
        halt (Table~\ref{tab:table1}).
  \item \textbf{Two deviations, both documented.} The assignment arm pose collides with the ToF sensor on this robot, so
        a hand-set pose 2\,cm higher is used, guarded by the planner. At that pose the camera cannot see the tag at
        20\,cm, so the camera brake is implemented but the stops come from the plan and odometry.
  \item \textbf{Possible improvements:} a lower ToF mount that allows the assignment pose (steeper camera, camera
        brake active), a larger tag for the 1.8\,m end of the range, and measured wheel speeds from the ESCs to compare
        with the modeled voltages.
\end{itemize}

\begin{notebox}
\small\textbf{\color{kfupmdark}\faFolderOpen\ Reproducing the results} (after \code{source ~/rm_setup.sh})
\begin{itemize}
  \item Build: \code{cd ~/ros2_ws && colcon build --symlink-install --packages-select rm_flatness_nav}.
  \item Run: driver in its own terminal, then \code{ros2 launch rm_flatness_nav mission.launch.py}; the run folder with
        both CSV files is printed at start.
  \item Offline tests: \code{cd ros2/rm_flatness_nav && python3 -m pytest test} (36 tests).
  \item Figures, generated tables and \code{analysis/results/summary.md}: \code{python3 analysis/run_all.py}.
\end{itemize}
\end{notebox}

%============================== REFERENCES ====================================
\clearpage
\isec{\faBook}{References}
\begingroup\small
\renewcommand{\section}[2]{}%
\begin{thebibliography}{99}\setlength{\itemsep}{3pt}
\bibitem{assign} A.~I. Bhatti, \emph{SCE 672 Advanced Robotics, Term Project Assessment 2: Differential Flatness
  Pursuit and Motor Voltage Dynamics}, KFUPM, 2026.
\bibitem{lectures} A.~I. Bhatti, \emph{SCE 672 lecture notes}, Lectures 5A (actuators and wheel kinematics),
  5B (polynomial trajectories) and 6A (differential flatness), KFUPM, 2026.
\bibitem{fliess} M.~Fliess, J.~L\'evine, P.~Martin, and P.~Rouchon, ``Flatness and defect of non-linear systems:
  introductory theory and examples,'' \emph{International Journal of Control}, vol.~61, no.~6, pp.~1327-1361, 1995.
\bibitem{murray} R.~M. Murray, M.~Rathinam, and W.~Sluis, ``Differential flatness of mechanical control systems: a
  catalog of prototype systems,'' in \emph{Proc. ASME International Mechanical Engineering Congress and Exposition}, 1995.
\bibitem{lynch} K.~M. Lynch and F.~C. Park, \emph{Modern Robotics: Mechanics, Planning, and Control}. Cambridge
  University Press, 2017, Ch.~9 and~13.
\bibitem{siegwart} R.~Siegwart, I.~R. Nourbakhsh, and D.~Scaramuzza, \emph{Introduction to Autonomous Mobile Robots},
  2nd ed. MIT Press, 2011.
\bibitem{olson} E.~Olson, ``AprilTag: A robust and flexible visual fiducial system,'' in \emph{Proc. IEEE International
  Conference on Robotics and Automation (ICRA)}, 2011, pp.~3400-3407.
\bibitem{wang} J.~Wang and E.~Olson, ``AprilTag 2: Efficient and robust fiducial detection,'' in \emph{Proc. IEEE/RSJ
  International Conference on Intelligent Robots and Systems (IROS)}, 2016, pp.~4193-4198.
\bibitem{pupil} Pupil Labs, \emph{pupil-apriltags: Python bindings for the AprilTag library},
  \url{https://github.com/pupil-labs/apriltags}.
\bibitem{corke} P.~Corke, \emph{Robotics, Vision and Control}, 2nd ed. Springer, 2017.
\bibitem{rmros} J.~Guzzi, \emph{robomaster\_ros: ROS 2 driver for DJI RoboMaster EP and S1},
  \url{https://github.com/jeguzzi/robomaster_ros}.
\bibitem{djisdk} DJI, \emph{RoboMaster SDK Documentation}, \url{https://robomaster-dev.readthedocs.io}.
\bibitem{ros2} S.~Macenski, T.~Foote, B.~Gerkey, C.~Lalancette, and W.~Woodall, ``Robot Operating System 2: Design,
  architecture, and uses in the wild,'' \emph{Science Robotics}, vol.~7, no.~66, 2022.
\bibitem{rep103} T.~Foote and M.~Purvis, \emph{REP 103: Standard Units of Measure and Coordinate Conventions},
  \url{https://www.ros.org/reps/rep-0103.html}.
\end{thebibliography}
\endgroup

\begin{notebox}
\small\textbf{\color{kfupmdark}\faLightbulb\ Acknowledgement of AI assistance.}
I used Claude (Anthropic) as an assistant in this assessment, at my direction: it wrote the code (the ROS 2 package
\code{rm_flatness_nav}, the analysis scripts and the helper tools) and drafted this report, explaining each part so
that I can present and defend it. I performed all hardware setup, every robot run and every measurement (tape
standoffs, tag size, range calibration), recorded the videos, and reviewed the code, the results and the text against
the raw logs in the repository.
\end{notebox}
